\documentclass[12pt,a4paper]{article}
\usepackage[utf8]{inputenc}
\usepackage[T1]{fontenc}
\usepackage[margin=2.5cm]{geometry}
\usepackage{hyperref}
\usepackage{enumitem}
\usepackage{tabularx}
\usepackage{booktabs}

\title{\textbf{EduRAG - Asynchronous AI Knowledge Base} \\ \large Milestone 0: Project Submission}
\author{Anna Ostrowska, Michał Kukla, Norbert Frydrysiak}
\date{\today}

\begin{document}

\maketitle

\section{Brief Description}
EduRAG is a highly scalable, AI-powered knowledge base utilizing the Retrieval-Augmented Generation (RAG) architecture. The platform allows users to upload massive volumes of textual data, such as gigabytes of textbooks, manuals, or research papers. To ensure a large-scale assumption, the system decouples the responsive user interface from the data processing layer. 

By leveraging a Serverless/PaaS event-driven architecture, uploaded documents are asynchronously read, chunked, and converted into embeddings in the background. Once processed, end-users can "chat" with the uploaded documents via a responsive UI, while the system dynamically retrieves the most relevant context to answer their queries.

\section{Main Features}
\begin{itemize}[label=\textbullet]
    \item \textbf{Asynchronous Document Ingestion Pipeline:} Direct-to-storage uploads trigger event queues. Serverless background workers consume these events to perform heavy PDF parsing, text chunking, and embedding generation without blocking the frontend.
    \item \textbf{Interactive Chat UI:} A fast, lightweight web interface where users can ask questions in natural language. The UI connects to an API that performs similarity searches and streams LLM-generated answers.
    \item \textbf{Scalable Vector Storage:} The generated text embeddings and metadata are stored in a Vector Database optimized for semantic search.
    \item \textbf{Real-Time Processing Status:} Users uploading large datasets can monitor the progress of the vectorization jobs via the dashboard.
    \item \textbf{Infrastructure as Code (IaC):} All cloud resources, including storage buckets, serverless functions, API gateways, message queues, and databases, are fully provisioned and managed using Terraform.
\end{itemize}

\section{System Architecture \& Implementation Strategy}
To fulfill the serverless and PaaS requirements, the system will be built using cloud services (e.g. AWS). The architecture is strictly divided into two decoupled workflows:

\subsection{Workflow 1: Asynchronous Ingestion Pipeline (Heavy Load)}
\begin{enumerate}
    \item \textbf{Upload:} The frontend requests a pre-signed URL from the API and uploads a large PDF, completely bypassing the backend servers.
    \item \textbf{Event Trigger:} A successful upload triggers a storage event, which drops a message into a Message Queue.
    \item \textbf{Background Processing:} A Serverless Function consumes the message, downloads the PDF, and uses a library (like LangChain) to split the text into smaller, overlapping chunks.
    \item \textbf{Vectorization \& Storage:} The function calls an Embedding API (e.g., OpenAI or HuggingFace) to convert text chunks into vectors, which are then stored in a Vector Database (e.g., ChromaDB).
\end{enumerate}

\subsection{Workflow 2: Real-time Inference Pipeline (Low Latency)}
\begin{enumerate}
    \item \textbf{Querying:} The user types a question in the UI. The request is routed via an API Gateway to a lightweight Serverless Function.
    \item \textbf{Similarity Search:} The function converts the user's question into a vector and queries the Vector DB for the top-K most semantically similar text chunks.
    \item \textbf{LLM Generation:} The retrieved chunks are injected into a prompt as context. The Large Language Model generates a final response based strictly on the provided context and streams it back to the UI.
\end{enumerate}

\section{Real-World Example}
To demonstrate the platform's utility, consider a university onboarding scenario:
\begin{itemize}[label=\textbullet]
    \item \textbf{Context Upload:} A professor uploads a 50-page course syllabus and grading criteria (PDF).
    \item \textbf{User Query:} A student asks the Chat UI: \textit{"What is the penalty for a late project submission?"}
    \item \textbf{RAG Action:} The Vector DB instantly retrieves the specific paragraph about delays. The LLM formulates the answer: \textit{"According to the syllabus, an M1 delay costs 5 points per week, while an M2 delay incurs a 10-point penalty."}
\end{itemize}

\section{Roles of Typical Users}
\begin{itemize}[label=\textbullet]
    \item \textbf{Knowledge Contributor (Admin/Teacher):} Has the authorization to upload large documents, manage the knowledge base, and monitor the status of background vectorization jobs.
    \item \textbf{End User (Student/Employee):} Interacts exclusively with the Chat UI. They can select a specific "knowledge context" (e.g., a specific textbook) and ask questions to retrieve summarized, AI-generated answers based solely on the uploaded files.
\end{itemize}

\section{Requirements Fulfillment}
\begin{itemize}[label=\textbullet]
    \item \textbf{UI:} A responsive web interface featuring a direct-to-cloud file upload portal with asynchronous progress tracking, alongside an interactive chat interface.
    \item \textbf{Storage:} Multi-tiered storage strategy utilizing Object Storage (for raw PDFs) and a managed Vector Database (for embeddings and fast semantic search).
    \item \textbf{Terraform:} All infrastructure (API Gateways, queues, serverless functions, storage buckets, and IAM roles) will be provisioned using Infrastructure as Code.
    \item \textbf{Large Scale Assumption:} Direct storage uploads prevent network bottlenecks. Decoupling the ingestion pipeline via message queues ensures that sudden spikes in large document uploads do not degrade the performance of the Chat UI.
    \item \textbf{PaaS / Serverless Preference:} The entire architecture relies on managed services (e.g., AWS).
\end{itemize}

\section{Team Roles \& Responsibilities}
To ensure every team member engages with cloud technologies, the workload is divided by \textbf{Cloud Domains} rather than traditional software layers. Every member is responsible for writing the Terraform configuration for their respective cloud resources.

\begin{enumerate}
    \item \textbf{Cloud Storage \& Edge Architect:} Responsible for the cloud entry points. Manages the deployment of the Chat UI via Cloud Hosting, configures Object Storage for secure, direct-to-cloud document uploads, sets up the API Gateway, and writes the associated Terraform scripts.
    \item \textbf{Asynchronous Compute \& Messaging Engineer:} Focuses on the event-driven core. Responsible for provisioning Message Queues, configuring Event Triggers, deploying the Serverless Background Workers that chunk the PDFs, and writing the associated Terraform scripts.
    \item \textbf{Cloud Data \& ML Integration Lead:} Manages data persistence and AI services. Responsible for provisioning the managed Vector Database, configuring NoSQL tables for background job metadata, integrating LLM APIs, and writing the associated Terraform scripts.
\end{enumerate}

\section{Project Schedule}
The project follows the official course milestones, with internal deadlines set to ensure timely delivery and testing.

\vspace{0.5cm}
\noindent
\begin{tabularx}{\textwidth}{@{} l l X @{}}
\toprule
\textbf{Date} & \textbf{Phase} & \textbf{Tasks \& Milestones} \\
\midrule
\textbf{Apr 17, 2026} & \textbf{M0} & \textbf{Project Submission:} Submit this document outlining project idea. \\
\addlinespace
Apr 18 -- 26 & Design & Draw microservices diagram, define API contracts (REST), and write initial Terraform for storage. \\
\addlinespace
\textbf{Apr 27, 2026} & \textbf{M1} & \textbf{Document \& Plan Presentation:} Present microservices, storage logic, and working Terraform. \\
\addlinespace
May 01 -- 22 & Build & Develop Serverless workers, integrate LLM API, connect Frontend, and test the event queue. \\
\addlinespace
May 23 -- Jun 02 & Scale Test & Perform load testing with massive PDFs to validate the "large scale assumption" and tune queues. \\
\addlinespace
\textbf{Jun 03, 2026} & \textbf{M2} & \textbf{Final Presentation:} Live demo of the asynchronous uploading and real-time chat processing. \\
\bottomrule
\end{tabularx}

\end{document}